\section{A wall is a floor turned on its side}\label{sec:wall}
The leverage list of the previous part opened with that sentence, and this section is what happened when it was built. The
sideways synthesis of section~\ref{sec:round7} could already place a part in any of the 24 orientations, but it chose those parts
the way the upright phases choose theirs --- one at a time, against the field, with a host invented for each --- and on a curved
flank it mostly declined. What was missing was not a candidate set. It was the \emph{frame}: on a vertical flank the stud axis of
the model is the wrong axis, and every phase the engine already has (the narrow aligned skin, the plate fill, the merges) is
written for a floor. Turn the field instead of the parts and all of them apply unchanged.

\subsection{The construction}
\code{motifs/wall.js}, the option \code{wall} (the \emph{Sideways building (SNOT)} switch), runs \textbf{first}, before the disc
layers and before the upright skin. That order is forced: after those two there is nothing left for it to take, because the rim
plates of a disc layer are exactly the staircase the sideways parts are there to replace (on the sphere, running it last leaves
173 candidates of which 166 report ``host busy''). Per facing it does five things.

\begin{enumerate}
\item \textbf{Resample.} The four horizontal facings $\pm x$, $\pm z$ are frames (\code{FACINGS}): local $+y$ is the facing, local
$+z$ is world $+y$, so a local \emph{row} is a world \emph{level} and a local \emph{level} is a slab of the model along the facing.
Each frame is one of the 24 orientation matrices of \code{orient.js}, which is what lets a piece be turned back exactly. The
remaining field is resampled in that frame: the 4\ldu{} samples along the facing pair into 8\ldu{} ``plates'', the plates along $y$
split into two samples.
\item \textbf{Solve.} A \code{Solver} is built over the resampled field and the ordinary phases run in it: the narrow aligned skin
(slopes, curved slopes, cheese, tiles, 1-wide, core parts only) and then a local fill of sideways plates and bricks that stacks
back onto the host face the way upright plates stack up.
\item \textbf{Turn back.} A local piece's yaw composed with the frame rotation is one of the 24 orientations; its box is the mapped
local box, with fractional $i$, $j$, $b$ in the world grid.
\item \textbf{Commit, bottom-up.} A piece is kept only if something holds it: a side-stud host brick in the cell behind its
underside, yawed so the stud faces it, or the studs of a sideways piece already committed beneath it. The host and everything on
it share an assembly id (\code{mi}) and are marked \code{rigid}, so \code{post.components} knows what holds what and no later pass
re-cuts them.
\item \textbf{Give up the cells.} The residual field in a committed piece's cells \emph{and in the level above} is zeroed. Without
the level above, the mop-up phases set 1$\times$1 plates down on sideways parts that have no stud under them --- 30 such orphans in
the first version, which the weld then paid for in plates.
\end{enumerate}

\subsection{Two lattices and a mask}
Where a sideways piece may sit is not a matter of taste; it is two exact lattices, and the local solver is told about them through
a new \code{Solver.allow} hook (figure~\ref{fig:walllattice}).

\textbf{Along the facing}, the underside of a sideways piece must land on a brick face --- a stud-cell boundary of the world grid.
The local level lattice is 8\ldu, so the whole pass runs twice with the levels offset by 0 and 4\ldu{} (\code{wallFoffs}); between
them they reach every face.

\textbf{Along $y$}, every LEGO side stud sits 14\ldu{} above the bottom of its brick. That number is read off the LDraw files
(\code{stud2a} 10\ldu{} below the top of 87087, 4070, 30414, 11211, 4733, 47905) and it agrees with every stud point the SNOT
synthesis of section~\ref{sec:round7} had recovered from the mined motifs (brackets: 30 of 40). A sideways part's own stud rows are
20\ldu{} apart, so only rows at $14 + 8k$ can meet a host at all, and with hosts restricted to brick layers --- \code{levelMod},
without which 26 of 33 hosts came out floating --- only rows at $14 + 24k$. Five offsets of the row lattice (\code{wallOffsets}, in
4\ldu{} samples) cover them: 4 hosts rows at $14 + 120k$, 3 at 38, 2 at 62, 1 at 86, 0 at 110.

A \textbf{mask} per facing says where it is allowed at all. A cell qualifies when it holds the \emph{silhouette} along the facing
(the outermost solid sample of its row at that level), when it is \emph{partial} (\code{wallPartial} $.12\ldots.88$) and when the
silhouette there is \emph{steep} (within \code{wallSteep} $60^\circ$ of vertical, read one level up and one down). The silhouette
rather than the top-surface gradient, because the gradient is blind exactly at the rim, where the neighbouring columns have no
surface at all. Partial, because a wall that already sits on the stud lattice is a brick's own face --- full or empty cells, nothing
to gain. And \emph{every} world cell a piece touches must be in the mask: a sideways part half inside a full cell would leave the
other half of that cell to upright parts that then cannot fit beside it.

\begin{figure}[h]
\centering
\begin{tikzpicture}[x=1cm,y=1cm,font=\scriptsize]
  % ---- left: the facing frame, seen from above
  \begin{scope}
    \fill[grey2] (0,0) -- (2.0,0) -- (2.45,1.5) -- (2.0,3.0) -- (0,3.0) -- cycle;
    \node[ink2,anchor=south west] at (0,3.25) {the model, seen from above};
    \node[ink2] at (0.95,1.9) {solid};
    \draw[->,blue1,thick] (2.2,1.5) -- (2.2,2.8) node[above,blue1]{local $+x$};
    \draw[->,orange1,thick] (2.2,1.5) -- (3.5,1.5);
    \node[orange1,anchor=south west,align=left] at (2.45,1.6) {local $+y$:\\the facing $f$};
    \node[ink2,anchor=north west,align=left] at (0,0.6) {local $+z$ = world $+y$, so a local row\\is a world level and every phase\\written for a floor runs unchanged};
  \end{scope}
  % ---- right: the elevation, the host and the two lattices (y in LDU)
  \begin{scope}[shift={(8.2,-0.55)},y=0.0155cm]
    \fill[grey2] (-0.8,0) rectangle (0,290); \draw[grey1] (-0.8,0) rectangle (0,290);
    \node[ink2,rotate=90] at (-0.4,215) {inside the model};
    \foreach \l in {0,8,...,288} \draw[grey2] (0,\l) -- (1.3,\l);        % world levels, 8 LDU
    % the host brick, 24 LDU, with its side stud 14 LDU up
    \fill[violet1!20] (-0.8,96) rectangle (0,120); \draw[violet1] (-0.8,96) rectangle (0,120);
    \fill[violet1] (0,110) rectangle (0.14,118);
    \draw[violet1] (-0.4,96) -- (-0.4,34); \node[violet1,anchor=north,align=center] at (-0.4,32) {a host brick,\\24\ldu, on a\\brick level};
    \draw[<->,violet1] (-1.0,96) -- (-1.0,114); \node[violet1,anchor=east] at (-1.05,105) {14\ldu};
    % the sideways pieces, their stud rows on the host's stud row and 20 LDU above it
    \fill[aqua1!30] (0.14,104) rectangle (1.15,124); \draw[aqua1] (0.14,104) rectangle (1.15,124);
    \fill[aqua1!30] (0.14,124) rectangle (0.85,144); \draw[aqua1] (0.14,124) rectangle (0.85,144);
    \draw[<->,aqua1] (1.28,114) -- (1.28,134); \node[aqua1,anchor=west] at (1.34,124) {20\ldu};
    \node[aqua1,anchor=west,align=left] at (1.34,225) {sideways pieces, stacking\\outward from the host face;\\their stud rows are 20\ldu{} apart};
    % the rows a host can offer
    \foreach \y in {14,38,62,86} \draw[orange1,thick] (-0.12,\y) -- (0.26,\y);
    \node[orange1,anchor=west,align=left] at (1.34,50) {hostable rows: $14 + 24k$};
    \draw[ink2,->] (-1.55,0) -- (-1.55,290) node[above,ink2]{world $y$};
  \end{scope}
\end{tikzpicture}
\caption{The frame and the two lattices. Left: per facing the field is resampled with the facing as ``up'', so the ordinary phases
--- written for a floor --- run on a wall unchanged. Right, in elevation: a side-stud host brick offers its stud 14\ldu{} above its
own bottom, and only brick levels may carry one, so the rows a sideways piece can hang from are $14 + 24k$; its own stud rows are
20\ldu{} apart, so it stacks \emph{outward} from the host face in 20\ldu{} steps, and the pass is repeated at five offsets of this
row lattice and two offsets of the level lattice along the facing.}
\label{fig:walllattice}
\end{figure}

\subsection{Hosts}
A host is a brick with studs on a vertical face. \code{HOSTS} in \code{wall.js} records their stud positions in the part's own
frame, read off the LDraw files, with \code{levelMod}: the 1$\times$1 brick \textbf{87087} (one stud, 14\ldu{} up, brick layers), and
the two bricks of height $1\frac{2}{3}$, \textbf{22885} (1$\times$2) and \textbf{32952} (1$\times$1), whose 40\ldu{} height \emph{is}
the sideways lattice --- they carry two stud rows, at 10 and 30\ldu, so a stack of them offers a side stud every two studs, and
every hostable row is reachable. 67329, 47905 and 4733 are in the table as well. The headlight brick 4070 is left out: its stud
face is recessed 4\ldu, so the part it holds would overlap its own cell.

The choice was made before the corpus of section~\ref{sec:mocs} could confirm it, and the confirmation is one of the pleasant
results of this work: the most frequent multi-model shaped motifs mined from 969 user-built models include \emph{stacked 22885
rows} (35 models, 304 uses), stacked 32952 (19 models), and ``22885 plus two sideways tiles'' --- that is, human builders build the
same host lattice, and the engine's sideways skin is a rediscovery of a known habit rather than an invention.

\subsection{Verifying it on the finished model}
A skin that wins on the raw solve can lose on the finished one, and the other way round, because the post-process is where most of
the pieces are decided. So the skin is checked the way the motifs are (section~\ref{sec:motifscore}): the \emph{whole finish} runs
with it and without it, and the job score $\mathrm{IoU} - 0.002\,\lvert\text{pieces}\rvert$ decides (\code{wallVerify}, on by
default; the grid phase is scored without the skin). On models whose flanks happen to land on the stud lattice --- the lying
cylinder, the rounded box --- the check simply drops it, which is the correct answer and not a failure.

\begin{table}[H]
\centering\small
\begin{tabular}{lrrrrrrl}
\toprule
 & \multicolumn{2}{c}{without the skin} & \multicolumn{2}{c}{with it} & \multicolumn{2}{c}{the skin itself} & \\
\cmidrule(lr){2-3}\cmidrule(lr){4-5}\cmidrule(lr){6-7}
model & pieces & IoU & pieces & IoU & pieces & hosts & the check \\
\midrule
teapot @24        & 583 & .7956 & \textbf{546} & \textbf{.8028} & 48 & 43 & kept \\
pony @24          & 618 & .6740 & \textbf{547} & .6713 & 52 & 39 & kept \\
sphere @16        & 766 & .8633 & \textbf{744} & .8552 & 58 & 50 & kept \\
dolphin @32       & 219 & .6704 & 221 & \textbf{.6813} & 21 & 14 & kept \\
\midrule
egg @16           & \textbf{967} & \textbf{.8554} & 972 & .8507 & 91 & 73 & dropped \\
hen @16           & \textbf{753} & \textbf{.8132} & 769 & .8099 & 64 & 58 & dropped \\
rounded box @24   & \textbf{616} & .9430 & 632 & .9435 & 10 & 6 & dropped \\
lying cylinder @16 & \multicolumn{2}{c}{186 p, .9144} & \multicolumn{4}{c}{never fires: its flank is on the stud lattice} & --- \\
\bottomrule
\end{tabular}
\caption{The sideways skin with every default of this document, both branches of its own verification. The teapot is the row to
read: \textbf{37 fewer pieces and more fidelity at the same time}, because a staircase of rim plates is replaced rather than
re-tiled; the pony pays .003 of IoU for 71 pieces, which the job score accepts. On the egg, the hen and the rounded box the skin
loses and is thrown away --- that is the check working, not the phase failing. The cost is 3--4$\times$ the solve time (40 local
solves: 4 facings $\times$ 2 level offsets $\times$ 5 row lattices) plus the second finish, which is why it is a switch and not a
default.}
\end{table}

\begin{figure}[H]
\centering
\begin{subfigure}{0.49\textwidth}\centering\includegraphics[width=\textwidth]{r_wall_teapot_off}\caption*{\small without: 583 pieces, IoU .796}\end{subfigure}
\begin{subfigure}{0.49\textwidth}\centering\includegraphics[width=\textwidth]{r_wall_teapot_on}\caption*{\small with: 546 pieces, IoU .803}\end{subfigure}
\caption{The teapot at 24 studs, rendered from the LDR export with real part geometry. Right: sideways pieces in aqua, the
side-stud host bricks that hold them in magenta (the 1$\times$2$\times$1$\frac{2}{3}$ bricks, their side studs visible). The skin is
a \emph{belt} around the steep part of the flank --- a Lowell sphere, not a SNOT hull --- because that is where the silhouette is
both steep and off the lattice. The top and the shoulder are left to the disc layers and the upright skin, which already do them
well.}
\label{fig:wallteapot}
\end{figure}

\subsection{What was measured and left out}
Two obvious extensions do not pay, and it is worth saying so.

\textbf{Wider hosts.} A row of 1$\times$1 hosts could be one 1$\times$2 (\textbf{11211}) or 1$\times$4 (\textbf{30414}) brick with
two or four side studs --- fewer pieces for the same studs. Both are now in \code{HOSTS} and reachable through \code{wallHosts}, and
on the eight benchmark models they are worth nothing. Behind 22885 they never fire at all: 22885 fails \emph{laterally}, in a
footprint a wider brick cannot fit either, and the 1$\times$1 hosts already cover the single-stud cases. Ahead of it the models
trade against each other inside $\pm.01$ of job score (the teapot and the dolphin gain, the sphere loses by as much). They stay
out of the default list.

\textbf{Brackets} (44728, 99780, 36840) are the hosts the corpus actually uses most after the headlight, and they need something
the engine cannot say yet: a piece whose body rises above --- or hangs below --- its own level. That is a change to the piece record,
not to this phase.

\section{What people build: the MOC corpus}\label{sec:mocs}
The motif library of section~\ref{sec:motifs} was mined from official sets. Official sets are designed by professionals under
constraints a hobbyist does not have (part budgets, moulds in production, a model that must survive a child), and the obvious
question is whether people building for themselves do something different. \code{docs/mocs/} holds 2,000+ Rebrickable MOC exports;
this section is the 1,069 of them under 1\,MB, 512\,MB of LDraw text, and what they changed.

\subsection{Reading a Studio export}
Three things had to be fixed before a single number was trustworthy, and all three are the difference between a file written by
LDraw's own tooling and one exported by BrickLink Studio, which is what people actually use.

\textbf{Embedded parts.} Studio writes every part it uses as its own \code{0 FILE 3023.dat} section inside the model, so a
100-part bench is 35\,KB of part geometry. The flattener was descending into those sections and reporting every primitive as an
unresolved sub-model; it now treats a section named after a library part as that part.

\textbf{A global offset.} The model sits where it was built: $-542.39$\ldu{} on $x$, $115.70$ on $z$ --- fractional LDU, not even a
stud multiple --- or 4\ldu{} up because it stands on a baseplate. Every part is then off the grid by the same residue.
\code{gridOffset} takes the modal 1\ldu{} residue of the upright catalogue parts' origins modulo a stud (in $x$, $z$) or a plate (in
$y$), uses the mean inside the modal bin as the offset (Studio writes fractional LDU, so 7.6 is 7.6, not 8), and subtracts it before
snapping. Modulo a \emph{whole} stud, not a half: measured modulo a half stud the registration settles on the jumper lattice and
16\,\% of all parts read as ``half'', where the true share is 3\,\%.

\textbf{A model built on its side.} A standing mosaic, or a file whose camera ``up'' became the model's up: all 2,953 parts
sideways, all off-grid. \code{modelUpright} counts the parts' up-vectors and, when more than half of the axis-aligned catalogue
parts agree on another axis, turns the whole model before registering.

Together these took the models with at least two grid pieces from 679 to 969 of 1,069, and the expressible share of the census from
47\,\% to 67\,\%. Neither touches how a model is \emph{solved}: they change what the census counts and what the miner sees.

\subsection{The census}
\begin{table}[H]
\centering\small
\begin{tabular}{lrrr}
\toprule
 & OMR (1,150 official) & 150 small MOCs & \textbf{1,069 MOCs $<$ 1\,MB} \\
\midrule
placements & 401,012 & 9,157 & \textbf{321,159} \\
expressible (placed + aliased) & 69.5\,\% & 69.5\,\% & \textbf{66.7\,\%} \\
sideways on the grid (SNOT) & 3.7\,\% & 10.9\,\% & \textbf{10.4\,\%} \\
sideways, off the grid & --- & 11.8\,\% & 5.6\,\% \\
tilted (hinges, clips on bars) & --- & 13.0\,\% & 10.4\,\% \\
off-grid after registration & --- & 9.7\,\% & 4.8\,\% \\
half-stud (jumpers, offset sub-assemblies) & --- & --- & 3.3\,\% \\
uncovered parts & 30.5\,\% & 15.7\,\% & 18.1\,\% \\
\bottomrule
\end{tabular}
\caption{The part census of \code{tools/omr\_stats.mjs} on the official sets and on the MOC corpus. The share of placements the
engine's model can express is the same; \emph{how} the rest is built is not.}
\end{table}

Two numbers carry the section. People build \textbf{three times as much sideways} as the official sets do --- 10.4\,\% of all
placements against 3.7\,\% --- and a tenth of everything they place is at an angle the stud grid cannot hold at all (hinges, clips
on bars, parts on pins). The first says the sideways skin of section~\ref{sec:wall} is aimed at the right thing and should
eventually be a default rather than a switch. The second is an honest ceiling on this whole method: a seventh of what a human
builder does is not a lattice problem and no catalogue will make it one.

\subsection{What the catalogue lacks, by family}
The 18\,\% uncovered placements sort into four families, and only one of them is fixed by measuring more parts.
\textbf{Connectors} (pins, clips, bars, handles, hinges: $\sim$9,000 uses) and \textbf{plants} are not shapes.
\textbf{Windows} --- a frame plus its glass filling a 1$\times$2$\times$2 or 1$\times$2$\times$3 box, 2,900 uses --- would matter for
buildings and for nothing organic; they are a volume part waiting to be made. \textbf{Tall solids and wedge plates} are shapes, and
14 of them were measured into the shape catalogue for this round (54 $\to$ 68 measured shapes, 266 $\to$ 280 parts): brick
1$\times$2$\times$5 (2454), 1$\times$2$\times$3 (22886), slope 75 2$\times$1$\times$3 (4460b) and 2$\times$2$\times$3 (3684a), slope
65 2$\times$2$\times$2 (3678b), and nine wings. The bulk-alias rule now reads a height suffix and the \code{\~{} \ldots{} (Obsolete)}
prefix, so ``\~{}Brick 1 x 2 x 5 without Centre Studs'' resolves to 2454 and brings 792 more placements with it. The brackets
(99780, 36840, 36841) were measured too and left out, for the reason given at the end of section~\ref{sec:wall}.

\subsection{Mining, merging, and whether it pays}
Mined with the same window miner as the official sets (3+ occurrences in 2+ models): 969 models, 246,909 pieces, 963,443 windows,
144,862 distinct groups, \textbf{9,223 motifs}, 222 of them sideways, and \textbf{6,226 that the official-set library does not
contain}. Merging is addition --- the same canonical key means the same assembly, so counts and model counts add up --- and gives
\textbf{20,313 motifs}.

What is genuinely new is less exciting than the count suggests, and the honest reading matters more than the number. The bulk of
the new multi-model motifs are \textbf{tilings}: grids of 2$\times$2 plates, rows of 2$\times$6, stacks of 2$\times$4 bricks, tile
floors --- the modular-building habit, and a mosaic contributes thousands of windows that are all the same. Of the shaped ones, the
informative group is the sideways lattice already quoted in section~\ref{sec:wall}, plus rows of inverted tiles, pairs of
round-end plates, and 2$\times$3 wings stacked in pairs as a two-plate-thick wedge.

\begin{table}[H]
\centering\small
\begin{tabular}{lrrrrr}
\toprule
8 models, 16 studs, verification \emph{off} & pieces & mean IoU & mean job score & from motifs & s \\
\midrule
motifs off & 10,988 & .7330 & $-2.014$ & 0 & 28 \\
the official-set library & 11,140 & .7334 & $-2.052$ & 2,435 & 54 \\
the MOC library alone & 11,401 & .7366 & $-2.114$ & 3,121 & 43 \\
merged & 11,320 & .7355 & $-2.095$ & 3,027 & 65 \\
\bottomrule
\end{tabular}
\caption{Every library buys IoU with pieces, and on the job score every one of them loses to ``motifs off'' on average. This is
the measurement that forced the verification of the next subsection; without it the right decision would have been to ship no
library at all.}
\end{table}

\subsection{Verification on the job score}\label{sec:motifscore}
The motif check of section~\ref{sec:motifs} compared the two \emph{raw} solves on IoU alone, with a margin of .01. Two things are
wrong with that. A library can win .003 of IoU for 4\,\% more pieces, which the job score --- and the user --- count as a loss. And
the comparison is made at the wrong moment: the raw solve \emph{with} assemblies has \emph{more} pieces than the one without
(rigid assemblies of small parts), while the finished one has fewer, because the merges can work on what the assemblies left. Ice
cream at 16 studs: 1,751 against 1,742 pieces raw, .5992 against .5965 IoU --- dropped --- but 1,131 against 1,157 pieces and .543
against .534 IoU once finished, which is a clear win.

So the check now wraps the whole finish, exactly like the sideways skin's: the model is finished with and without the assemblies
and the one that wins on $\mathrm{IoU} - 0.002\,\lvert\text{pieces}\rvert$ is returned (\code{motifScore}, on by default;
\code{motifGain} is the margin on that score, 0). \code{post.motifCheck} reports both IoUs and both piece counts.

\begin{table}[H]
\centering\small
\begin{tabular}{lrrrl}
\toprule
8 models, 16 studs, defaults & pieces & mean IoU & mean job score & kept on \\
\midrule
motifs off & 10,988 & .7330 & $-2.014$ & --- \\
merged library, verified & \textbf{10,877} & \textbf{.7358} & \textbf{$-1.984$} & delfin, felix, ice cream, table \\
\bottomrule
\end{tabular}
\caption{Never worse by construction, better on half the models. At 24 studs: 29,593 pieces / IoU .7999 without motifs against
29,445 / .8030 with the verified library --- fewer pieces \emph{and} more fidelity.}
\end{table}

\section{Making it run in a browser}\label{sec:fast}
The merged library doubled the data the engine carries, and the engine runs in a Web Worker pool inside a browser tab, eight
copies of it, next to a three.js viewport holding the model. This section is the round that made that affordable. It is
engineering rather than method, but two of the five changes are exact statements about the solver and belong in the record.

\subsection{Data that no run reads}
\code{motifs.js} was 8.09\,MB of object literals that \emph{every worker} parsed on import, whether or not motifs were switched on.
The observation that fixes it is that the mined \textbf{key already encodes every part}: \code{w,d,h|id:rot:i,j,b \ldots}, with
offsets in tenths of a stud and half plates. The \code{parts} array of the JSON was pure duplication. The library now ships as one
line per motif, \code{<key>\textbackslash t<n>\textbackslash t<models>}, with everything no setting can reach left out (fewer than 3
occurrences, fewer than 2 models, more than 16 parts, a part the catalogue does not have): \textbf{10,669 of the 20,313 motifs,
0.83\,MB instead of 8.09}. \code{parseLibrary} rebuilds the records on first use in 44\,ms, so a run with motifs off never touches
them. Round-tripped against the JSON: the same parts for all 10,669 motifs, the same 1,942 motifs kept, the same 8,724 variants in
the same order, the same finished models piece for piece.

The measured part catalogues went the same way the next morning. They are geometry, not a key, so the trick is different: the
height fields are nearly constant (a plain brick is 24\ldu{} everywhere and 0 underneath), so each part is one line --- the small
fields as JSON, the \code{top} / \code{bot} / \code{cover} maps run-length encoded at 0.01\ldu{} --- decoded on demand into flat
\code{Float32Array}s. Together: \textbf{3.94\,MB $\to$ 0.18\,MB}. The drawing geometry (2.2\,MB of triangles) moved into a module of
its own that only the exporter imports, so a solver worker never loads it at all and a part's triangles are decoded the first time
it is drawn. The quantisation step is not free to choose: at 0.1\ldu{} two of the five benchmark models shifted by four or five
pieces, a tie in the scoring flipping, so it is 0.01.

\begin{figure}[h]
\centering
\includegraphics[width=\textwidth]{moc_library}
\caption{Left: the shipped library and what the default settings do with it. Four fifths of the file --- the mosaic tilings above
all --- is rejected at run time by \code{motifShapedOnly} and the part-count cap, which raises the obvious question of cutting them
at pack time. Measured: it would save 0.57\,MB and \emph{nothing} in time (the rejection happens before any volume is built: 1,811
against 1,798\,ms to build the library), while removing a knob the user can turn. Not done. Right: the library as mined, merged,
and as it ships.}
\label{fig:moclibrary}
\end{figure}

\subsection{The candidate scan}
Where the motif phase's time actually went, counted on the hen at 16 studs: 0.77\,M positions survive the box-sum gate, 4.1\,M
(position, variant) pairs pass the cheap volume tests, 0.83\,M reach the exact per-cell test --- and \textbf{628\,M cells}, because
that test summed every cell of the variant before deciding, while 92\,\% of those tests fail. Two exact changes in
\code{Solver.candidates}, neither of which changes a single placement:

\textbf{An admissible window.} Inside a box group both cheap tests are monotone in the variant's own volume ($s_w \ge
\texttt{min\_cov}\cdot v_{\text{tot}}$ and $\lvert s_w - v_{\text{tot}}\rvert \le \texttt{maxErr}$), so for a given position the
variants that can pass are a \emph{contiguous run} of the group once it is sorted by volume. A binary search for the two ends
replaces the walk over all of them --- a motif group holds up to 362.

\textbf{An early exit.} The accumulated error \code{over} only grows, and a candidate is accepted only if it finishes below the
acceptance cap (and below what a flat stack of plates would cost, where \code{beatFlat} applies). So the cell loop stops the moment
it has already failed: \textbf{628\,M cells $\to$ 124\,M}.

Both help every phase, not only the motifs, and the models are identical piece for piece (hen, dolphin, table, teapot, pony:
identical piece multisets, identical IoU). A third change in the same spirit: \code{assemblyVariants} deduplicated rotations by
building \code{V.map(x => x.toFixed(4)).join(',')}, a $\sim$16\,KB string per variant for 8,724 variants; two 32-bit FNV walks over
the cells quantised to $1/4096$ take the library build from 2,328 to 672\,ms.

\subsection{Finishing a model three times instead of four}
With both verifications on --- the sideways skin's and the assemblies' --- the two checks nested: every combination of
\{skin\}~$\times$~\{assemblies\} was finished, and the expensive corner, the skin \emph{and} the motif phase, ran twice. The
assemblies are now decided on the branch \emph{without} the skin, which is the cheap one and is also the skin check's own
``without'', and the skin is then finished once with that decision. Three finishes. Dolphin at 16 studs with both: 19.8\,s $\to$
11.9\,s; pony 27.7 $\to$ 18.6; identical models and identical decisions on both. The colour reference is memoised per mesh at the
same time (the 26-direction visibility buffer and the point grid depend on the mesh, not on the pieces): 345\,ms $\to$ 74 on the
hen's second finish, and the saving grows with the mesh.

\subsection{The worker pool}
The last one is a plain bug with a large effect, and the symptom was the one users report as ``the app gets slower the longer I
use it''. \code{runMethod} began by cancelling: it terminated the whole pool and spawned eight \emph{new} workers for every run,
each re-importing the engine from scratch, while the old copies waited for the garbage collector. The pool is kept now --- workers
are terminated only when a run is cancelled, when the pool shrinks, or when one is stale --- and the correctness that buys it is
cheap: every worker answer carries the run's key and a mismatched answer is dropped, so a late reply from the previous run cannot
land in the new one. After a run the idle workers are told to release the prepared field (megabytes) while keeping the modules.

\begin{figure}[h]
\centering
\includegraphics[width=\textwidth]{round13_fast}
\caption{What the three rounds moved. The engine import is what every one of the eight workers paid on every run before the pool
was kept; the heap is per worker. The last panel is the whole benchmark: eight models at 16 studs with motifs on and verified, at
identical output (10,877 pieces, mean IoU .7358).}
\label{fig:fast}
\end{figure}

\section{Buildability, and the booklet}\label{sec:booklet}
\subsection{Overlaps as a hard test}
``Pieces that intersect in space'' is the one failure a LEGO model cannot argue with, and until this round nothing checked for it.
\code{post.audit} rasterises every piece afresh at the engine's 4\ldu{} resolution --- a sideways piece by the whole cells its
fractional box touches --- and counts samples claimed twice; it runs at the end of every solve and is reported with the result. It
found two real bugs within a day of being written.

The first was the flatten fallback of the bridge pass. It replaced a smooth piece by a brick of ``the same footprint'' through
\code{mk(brick, brick.w === p.w ? p.w : p.d, \ldots)}, which \emph{transposes} the footprint whenever the brick's canonical $w$ and
$d$ are the other way round --- a 3$\times$1 curved slope became a 1$\times$3 brick lying across its neighbours --- and it never
looked above the first level. Ice cream at 16 studs: 175 overlapping samples $\to$ 85 after the footprint fix $\to$ 0 once the
flatten also checks the three levels of a brick.

The second is the more interesting one, because it is the kind of bug a fractional lattice invites. \code{post.bridge} built its
occupancy grid with \code{for (let l = p.b; l < p.b + p.h; l++)}, which silently does \emph{nothing} for a sideways piece sitting at
$b = 52.5$. The chain search therefore saw the entire sideways skin as empty air and routed plates straight through it: 790
overlapping samples in 47 cells on the pony, 50 on the dolphin. Both that pass and the exterior-air mask now go through one
function, \code{post.cellBox(p)} --- every whole cell a fractional box touches --- which \code{occupancy} and \code{bracing} already
used. Zero overlapping samples everywhere again, and it is a small \emph{gain} rather than a cost: dolphin 729 $\to$ 726 pieces at
$+.0013$ IoU, pony 2,437 $\to$ 2,419 at $+.001$. The rule for the rest of the codebase is worth stating once: \emph{a piece list is
not a grid of integers}, and any loop of that shape is a bug.

\subsection{Instructions}
A piece list is not a buildable model until somebody can build it, so the Studio exports a booklet shaped like a real set's.

\textbf{Steps} are the build order, not the layers. Pieces are grouped into levels and the levels into steps of at most a few
pieces, but a sideways piece is placed at the level of the highest \emph{upright} piece of its assembly, and within a step the host
comes before whatever hangs on it. Before that rule, 47 of the pony's 192 sideways pieces were drawn before anything held them
(the dolphin: 58); now none are.

\textbf{Each step is an isometric orthographic capture} of the model as built so far, framed on the step's own pieces with a
margin and never below about half the model, so the eye has somewhere to stand. The new pieces are in their own colours with black
outlines, everything already built is grey with grey outlines --- the one convention that makes an instruction page readable --- and
the call-out above the picture aggregates by part (``2$\times$4 plate $\times$4'') with a swatch per colour.

\textbf{The cover} is one path-traced view at the resolution the user chose in the app, the camera pulled in to the bounding sphere
with a 1\,\% margin and the image laid out \code{contain} over the whole page, with the logo, the title, a stable five-digit set
number, the warning box and the model's numbers above it. The page is \textbf{portrait when the render is taller than wide} and
landscape otherwise, decided from the PNG header, while the inside of the booklet is always landscape --- which is visible in the
printed PDF: page 1 is 210$\times$297\,mm and the rest 297$\times$210. A \emph{finished} page closes the build with two quick
viewport renders (no path tracer) framed on the model's real world box, and a parts list ends it.

Shooting a booklet is sixty-odd captures and a path-traced cover, so the export is cached: a signature over the piece list and the
settings that shape the pages returns the previous file untouched when nothing has changed.

\begin{figure}[p]
\centering
\begin{subfigure}{0.33\textwidth}\centering\includegraphics[width=\textwidth]{bk_cover}\end{subfigure}
\begin{subfigure}{0.64\textwidth}\centering\includegraphics[width=\textwidth]{bk_steps}\end{subfigure}\\[6pt]
\begin{subfigure}{0.49\textwidth}\centering\includegraphics[width=\textwidth]{bk_final}\end{subfigure}
\begin{subfigure}{0.49\textwidth}\centering\includegraphics[width=\textwidth]{bk_parts}\end{subfigure}
\caption{The booklet of the hen at 16 studs: 782 pieces, 31 different parts, 63 steps over 19 pages. The cover (portrait, because
this render is taller than it is wide), a step page, the finished-model page and the parts list. Everything here is produced from
the piece list alone, headlessly, by the same code the Studio button runs.}
\label{fig:booklet}
\end{figure}
